\BourbakiExercises{习题}

\begin{enumerate}
\def\labelenumi{\arabic{enumi})}
\tightlist
\item
  设 D 是一个体，V 是 D 上的一个右向量空间，A 是环 \(\operatorname{End}_{\mathbb{D}}(\mathrm{V})\)。对 V 的任一线性子空间 W，用 \(\mathfrak{a}(\mathrm{W})\)（相应地，b(W)）表示 A 中满足 aW = 0（相应地，aV ⊂ W）的元素 a 所成的集。
\end{enumerate}

\begin{enumerate}
\def\labelenumi{\alph{enumi})}
\item
  证明映射 \(W \mapsto \mathfrak{a}(W)\)（相应地，\(W \mapsto \mathfrak{b}(W)\)）是从 V 的线性子空间之集到 A 的由幂等元生成的左（相应地，右）理想之集上的一个递减（相应地，递增）双射。
\item
  A 的极小左理想是诸理想 \(\mathfrak{a}(\mathrm{H})\)，其中 H 是 V 中一个超平面。A 的极小右理想是诸理想 \(\mathfrak{b}(\mathrm{D})\)，其中 D 是 V 中一条直线。由此推出 A 的有限长度的左（相应地，右）理想的形式。
\item
  设 W 与 \(W'\) 是 V 的两个子空间，满足 \(W \neq 0\) 且 \(W' \neq V\)。证明交 \(\mathfrak{a}(W) \cap \mathfrak{b}(W')\) 不约化为 0。
\end{enumerate}

\(d\) ) 映射 \(u\mapsto{}^{t}u\) 定义了从 A 到 \(\mathrm{End}_{\mathrm{D}^{\circ}}(\mathrm{V}^{*})\) 的一个子环 B 的环同构。证明在此同构下，\(\mathfrak{a}(\mathrm{W})\)（相应地，\(\mathfrak{b}(\mathrm{W}))\)）的像是理想 \(\mathfrak{b}(\mathrm{W}^{\circ})\)（相应地，\(\mathfrak{a}(\mathrm{W}^{\circ}))\)）在 B 上的迹，其中 \(\mathrm{W}^{\circ}\) 表示 W 在 \(\mathrm{V}^*\) 中的正交

\begin{enumerate}
\def\labelenumi{\alph{enumi})}
\setcounter{enumi}{4}
\tightlist
\item
  设 \(\mathfrak{F}\) 是 V 的线性子空间所成的集。对 A 的任一左理想 \(\mathfrak{a}\)，用 \(\mathfrak{F}(\mathfrak{a})\) 表示 \(\mathfrak{F}\) 中由诸子空间 \(\operatorname{Ker} u\)（\(u\) 取遍 \(\mathfrak{a}\)）组成的子集。集 \(\mathfrak{F}(\mathfrak{a})\) 对交运算稳定，且 V 的每个包含 \(\mathfrak{F}(\mathfrak{a})\) 中一个元素的线性子空间都属于 \(\mathfrak{F}(\mathfrak{a})\)。证明映射 \(\mathfrak{a} \mapsto \mathfrak{F}(\mathfrak{a})\) 是从 A 的左理想之集到 \(\mathfrak{F}\) 的具有这两条性质的子集之集上的一个双射。f) 设 V 的维数无限，并用 \(\mathfrak{c}\) 表示它。设 \(\mathfrak{m}\) 是 A 的一个极大左理想，它不具有 \(\mathfrak{a}(D)\) 的形式。证明诸子空间 \(W \in \mathfrak{F}(\mathfrak{m})\) 的交约化为 0，这些子空间中的每一个都是无限维的，并且存在余维数为 \(\mathfrak{c}\) 的这样的子空间。
\end{enumerate}

¶ 2) 设 D 是一个体，V 是 D 上一个无限维 c 的右向量空间，A 是环 \(\operatorname{End}_{\mathbb{D}}(\mathbb{V})\)。

\begin{enumerate}
\def\labelenumi{\alph{enumi})}
\item
  设 u 与 v 是 V 的自同态，满足 \(\operatorname{rg}(u) \leqslant \operatorname{rg}(v)\)。证明存在 A 的元素 a 与 b，使得 u = avb。
\item
  对任一无限基数 \(n \leqslant c\)，把 V 的秩 \textless{} n 的自同态所成的集记作 \(A_{n}\)。证明 \(A_{n}\) 是 A 的一个双边理想，并且 A 的每个异于 \(\{0\}\) 与 A 的双边理想都是某个 \(A_{n}\)（利用 a)）。特别地，有限秩自同态的理想是 A 的最小非零双边理想，而理想 \(A_{c} = T\)（由秩 \textless{} c 的自同态组成）是 A 的异于 A 的最大双边理想。
\item
  设 B 是 V 的一个基，U 是 B 上的一个超滤子，它比由 B 的基数 \textless{} c 的子集的补集组成的滤子 G 更细。对 U 中每个集 U，用 \(e_{U}\) 表示 V 的投影算子，其核由 U 的向量生成，其像由 B-U 的向量生成。用 \(a_{U}\) 表示由诸 \(e_{U}\)（其中 U 取遍 U）生成的 A 的左理想。证明理想 \(T + a_{U}\) 异于 A。设 V 是 B 上另一个包含 G 且异于 U 的超滤子。证明我们有 \(a_{U} + a_{V} = A\)。d) 证明 B 上包含 G 的超滤子所成的集的基数是 \(2^{2^{c}}\)。（利用一般拓扑，I, §8, 第 135 页习题 6 与一般拓扑，I, §4, 第 126 页习题 5 的论证。注意在后一习题中，X 的每个非空开子集在 F 上的迹与 F 具有相同的基数。）
\item
  设 \(\operatorname{Card}(\mathrm{D}) \leqslant 2^{\mathfrak{c}}\)。由 \(d\) 推出，单 \((\mathrm{A} / \mathfrak{T})\)-模的类所成的集具有基数 \(2^{2^{\mathfrak{c}}}\)（参见 VIII, 第 52 页，习题 5, \(d\) )）。
\end{enumerate}

\(f\) ) 由上述推出，环 A 是本原的但不是拟单的，而环 \(\mathrm{A} / \mathfrak{T}\) 是拟单的但不是单的（VIII, 第 120 页，注 2）。

\begin{enumerate}
\def\labelenumi{\arabic{enumi})}
\setcounter{enumi}{2}
\tightlist
\item
  设 \(D\) 是一个体，\(V\) 是 \(D\) 上的一个右向量空间，\(A\) 是环 \(\operatorname{End}_D(V)\) 的一个子环。
\end{enumerate}

\begin{enumerate}
\def\labelenumi{\alph{enumi})}
\tightlist
\item
  证明下列条件等价：
\end{enumerate}

\begin{enumerate}
\def\labelenumi{(\roman{enumi})}
\item
  对每个自同态 \(u\)（\(V\) 的）与向量的每个有限序列 \(x_{1},\ldots ,x_{n}\)（在 \(V\) 中），存在一个元素 \(a\)（\(A\) 的），满足 \(u(x_{i}) = a(x_{i})\)（对 \(1\leqslant i\leqslant n\)）。
\item
  对每对元素 \(x, y\)（\(V\) 的，\(x \neq 0\)），存在 \(a \in A\)，使得 \(a(x) = y\)。
\end{enumerate}

当这些条件满足时，我们称 \(\operatorname{End}_{\mathbb{D}}(\mathrm{V})\) 的子环 A 是稠密的。从现在起设 A 是稠密的。

\begin{enumerate}
\def\labelenumi{\alph{enumi})}
\setcounter{enumi}{1}
\item
  证明 A 是本原的（VIII, 第 120 页，注 2）。反之，证明每个本原环都同构于某个向量空间的自同态环的一个稠密子环。
\item
  证明 A 的中心的非零元素都是单自同态。特别地，本原环的中心是一个整环。
\item
  设 V 是无限维的。对每个整数 n \textgreater{} 0，证明存在 A 的一个子环 \(A_{n}\) 和一个满同态 \(\mathrm{A}_{n} \to \mathbf{M}_{n}(\mathrm{D})\)。
\end{enumerate}

\begin{enumerate}
\def\labelenumi{\arabic{enumi})}
\setcounter{enumi}{3}
\tightlist
\item
  设 \(\mathrm{A}\) 是一个本原环。
\end{enumerate}

\begin{enumerate}
\def\labelenumi{\alph{enumi})}
\item
  若环 A 是交换的（相应地，左阿廷的），则它是一个体（相应地，一个单环）。
\item
  设 a 与 b 是 A 的非零左理想。证明理想 ab 不为零。\(^{(3)}\) 由此推出 A 的有限多个非零双边理想的交不约化为 0。

  \( ^{(3)} \) 在本习题与下一习题中，若 a 与 b 是 A 的左理想，则用 ab 表示由诸乘积 ab（\(a \in a\)，\(b \in b\)）生成的左理想。
\item
  证明若对每个 \(a \in A\)，元素 \(1 + a^{2}\) 在 A 中可逆，则 A 是一个体（利用习题 3, b)）。
\end{enumerate}

\begin{enumerate}
\def\labelenumi{\arabic{enumi})}
\setcounter{enumi}{4}
\tightlist
\item
  设 \(\mathrm{A}\) 是一个环。
\end{enumerate}

\begin{enumerate}
\def\labelenumi{\alph{enumi})}
\item
  设 e 是 A 中的幂等元。证明若 A 是本原的（相应地，拟单的），则环 eAe 亦然（利用习题 3, b)，VIII, 第 113 页的引理 4，以及 VIII, 第 15 页的习题 6, a)）。
\item
  设 n 是一个整数 \(\geqslant 1\)。证明矩阵环 \(\mathbf{M}_{n}(\mathrm{A})\) 是本原的（相应地，拟单的）当且仅当 A 是本原的（相应地，拟单的）（参见 VIII, 第 114 页，例 2）。
\item
  设 B 是一个与 A 森田等价的环；则 B 是本原的（相应地，拟单的）当且仅当 A 亦然。
\end{enumerate}

\begin{enumerate}
\def\labelenumi{\arabic{enumi})}
\setcounter{enumi}{5}
\tightlist
\item
  设 \(A\) 是一个拟单环。
\end{enumerate}

\begin{enumerate}
\def\labelenumi{\alph{enumi})}
\item
  一个 A-模是生成模当且仅当它的对偶非零。
\item
  设 \(\mathfrak{a}\) 是 A 的一个左理想，D 是环 \(\operatorname{End}_{\mathrm{A}}(\mathfrak{a})\)。证明 \(\mathfrak{a}\) 是一个有限生成的投射右 D-模，且映射 \(a \mapsto a_{\mathfrak{a}}\)（从 A 到 \(\operatorname{End}_{\mathrm{D}}(\mathfrak{a})\)）是双射。环 D 是拟单的当且仅当 A-模 \(\mathfrak{a}\) 是投射的且有限生成的。
\item
  证明一个包含极小左理想的拟单环是单的。
\end{enumerate}

\begin{enumerate}
\def\labelenumi{\arabic{enumi})}
\setcounter{enumi}{6}
\tightlist
\item
  设 A 是一个环。
\end{enumerate}

\begin{enumerate}
\def\labelenumi{\alph{enumi})}
\item
  设 \(\mathfrak{a}\) 是 A 的一个极小左理想。证明我们有 \(\mathfrak{a}^2 = 0\) 或 \(\mathfrak{a}^2 = \mathfrak{a}\)。在后一情形，证明存在一个幂等元 \(e\)（在 \(\mathfrak{a}\) 中），使得 \(\mathfrak{a} = \mathrm{A}e\)（若 \(a \in \mathrm{A}\) 满足 \(\mathfrak{a}a = \mathfrak{a}\)，考虑自同构 \(x \mapsto xa\)（\(\mathfrak{a}\) 的））。
\item
  设 e 是 A 中的幂等元。理想 Ae 是极小的当且仅当环 eAe 是一个体。
\item
  设 \(\mathfrak{a}\) 与 \(\mathfrak{b}\) 是 A 的两个同构的左理想。证明我们有 \(\mathfrak{ab} = \mathfrak{b}\)（若 \(\mathfrak{a}^2 = \mathfrak{a}\)），而 \(\mathfrak{ab} = 0\)（若 \(\mathfrak{a}^2 = 0\)）。
\end{enumerate}

¶ 8) 设 A 是一个环。A 的左底座（或简称底座）s 是 A-模 \(A_{s}\) 的底座，即（VIII, 第 65 页）A 的极小左理想之和。它是一个半单 A-模，而它的同型分量称为它的足。设 a 是 s 的一个足。

\begin{enumerate}
\def\labelenumi{\alph{enumi})}
\item
  证明 \(\mathfrak{a}\) 是一个双边理想，并且含于 \(\mathfrak{a}\) 的诸平方为零的极小理想（习题 7）之和是一个双边理想，即 \(\mathfrak{a}\) 与其零化子之交。
\item
  设 a 包含一个平方非零的极小理想。证明不存在 a 的非零元素 x 满足 ax = 0（利用习题 7, c)）。
\item
  在 b) 的假设下，证明伪环 a（I, §8, No.~1, 第 98 页）的每个左理想都是 A 的理想（注意 a 的极小理想就是 A 的含于 a 的极小理想）。
\item
  设 A 的底座 s 不包含任何平方为零的极小左理想。证明 A 的每个作为 A-模具有有限长度的左理想 b 都含于 \(\mathfrak{s}\)（对 \(\mathfrak{b}\) 的长度用归纳法，注意若 \(e\) 是 \(\mathfrak{b}\) 中的幂等元，则 \(\mathfrak{be}\) 是 \(\mathfrak{b}\) 的一个直因子）。
\end{enumerate}

¶ 9) 设 D 是一个体，V 是 D 上的一个右向量空间，A 是环 \(\operatorname{End}_{\mathrm{D}}(\mathrm{V})\) 的一个稠密子环（习题 3）。

\begin{enumerate}
\def\labelenumi{\alph{enumi})}
\item
  子环 A 包含一个非零的有限秩自同态，当且仅当 A 的底座 s 非零（利用习题 4, a) 与 7, b) 证明 A 的极小理想由一个幂等元生成，再证明该幂等元的秩为 1）。此时 A 的底座只有单个足（习题 4, a)）。
\item
  从现在起设 \(s \neq 0\)。设 a 是 A 的一个极小左理想。证明存在一个元素 \(\ell\)（V 的对偶 \(V^{*}\) 的），使得 a 由诸自同态 \(z \mapsto v \ell(z)\)（\(v \in V\)）组成。设 \(V'\) 是由线性型 \(\ell \in V^{*}\) 组成的集，这些线性型使得自同态 \(z \mapsto v \ell(z)\) 对每个 \(v \in V\) 都属于 A。证明 \(V'\) 是 \(V^{*}\) 的一个 (D, A)-子双模，且 s 是 V 的这样的有限秩自同态组成的集：其核可写成 \(\operatorname{Ker} \ell_{1} \cap \cdots \cap \operatorname{Ker} \ell_{n}\)，其中 \(\ell_{1}, \ldots, \ell_{n}\) 属于 \(V'\)。仅当向量空间 V 是有限维且 \(A = \operatorname{End}_{\mathrm{D}}(\mathrm{V})\) 时，才能有 s = A。
\item
  设 \(\mathfrak{b}\) 是 A 的一个极小右理想。证明存在向量 \(v \in V\)，使得 \(\mathfrak{b}\) 由诸自同态 \(z \mapsto v\ell(z)\)（\(\ell \in V'\)）组成。由此推出 A 的右底座等于其左底座 \(\mathfrak{s}\)。
\end{enumerate}

\(d\) ) 证明 \(\mathfrak{s}\) 还是 A 的最小非零双边理想。

\begin{enumerate}
\def\labelenumi{\alph{enumi})}
\setcounter{enumi}{4}
\tightlist
\item
  由 b) 推出，一个具有非零底座的（左或右）诺特本原环是单的（通过考虑含于 s 中的理想证明 V 必为有限维）。
\end{enumerate}

\(f\) ) 设 M 是 V 的一个有限维子空间。证明 \(\mathfrak{s}\) 包含 V 的一个以 M 为像的投影算子 \(e_{\mathrm{M}}\)。（设 \(\mathrm{N}^{\prime}\subset \mathrm{V}^{\prime}\) 是 M 的正交的一个补子空间。证明存在 M 的基 \((v_{j})_{1\leqslant j\leqslant k}\) 与 \((\ell_i)_{1\leqslant i\leqslant k}\)（\(\mathrm{N}'\) 的基），使得 \(\ell_i(v_j) = \delta_{ij}\)，并考虑自同态 \(x\mapsto \sum_{i}v_{i}\ell_{i}(x).)\) 设 N 是 \(e_{\mathrm{M}}\) 的核。证明 V 的每个自同态 \(u\)（满足 \(u(\mathrm{M})\subset \mathrm{M}\) 且 \(u(\mathrm{N}) = 0\)）都属于 \(\mathfrak{s}\)。

\(g)\) 设 \(J\) 是 \(A\) 的一个有限子集。证明存在一个分解 \(V = M \oplus N\)，其中 \(M\) 是一个有限维子空间，\(N\) 是 \(V'\) 的一个有限子集的正交，使得自同态 \(u\)（\(V\) 的）（满足 \(u(M) \subset M\) 且 \(u(N) = 0\)）组成的环含于 \(A\) 且包含 \(J\)。（设 \(M_1 = \sum_{u \in J} \operatorname{Im} u\) 与 \(N_1 = \cap_{u \in J} \operatorname{Ker} u\)，把 \(f\) ) 应用于子空间 \(M\)（包含 \(M_1\) 且使 \(M + N_1 = V\)），然后取 \(N\)（含于 \(N_1\) 中）。）

\begin{enumerate}
\def\labelenumi{\arabic{enumi})}
\setcounter{enumi}{9}
\item
  设 A 是一个具有非零底座 \(\mathfrak{s}\) 的本原环，M 是一个 A-模。则 M 是忠实的且同型的，当且仅当我们有 \(\mathfrak{sM} = M\)（利用 VIII, 第 50 页的命题 4）。由此推出，若 \(0 \to M' \to M \to M'' \to 0\) 是 A-模的一个正合列，且 \(M'\) 与 \(M''\) 都是忠实的且同型的，则 M 亦然。
\item
  设 \(\mathrm{K}\) 是一个交换域，\(\mathrm{V}\) 是 \(\mathrm{K}\) 上的一个以 \((e_n)_{n\geqslant 1}\) 为基的向量空间。设 \(u\) 与 \(v\) 是 \(\mathrm{V}\) 的自同态，由 \(u(e_1) = 0\)、\(u(e_p) = e_{p - 1}\)（对 \(p > 1\)）与 \(v(e_p) = e_{p + 1}\)（对 \(p\geqslant 1\)）定义，并设 \(\mathrm{A}\) 是 \(\operatorname {End}_{\mathrm{K}}(\mathrm{V})\) 的由 \(1, u,\) 与 \(v\) 生成的 K-子代数。我们有 \(uv = 1_{\mathrm{V}}\) 且 \(vu\neq 1_{\mathrm{V}}\)。
\end{enumerate}

\begin{enumerate}
\def\labelenumi{\alph{enumi})}
\item
  证明诸元素 \(v^i u^j\)（\(i \geqslant 0, j \geqslant 0\)）构成 A 在 K 上的一个基。
\item
  证明环 A 是本原的，且其底座 s 是 V 的这样的自同态 w 组成的集：序列 \((w(e_{p}))_{p\geqslant1}\) 具有有限支集（注意 s 是 A 的包含 \(1_{V}-vu\) 的最小双边理想）。证明 K-代数 A/s 同构于 \(K[X,X^{-1}]\)。
\end{enumerate}

¶ 12) 设 A 是一个环，d 是 A 的一个导子。把环 \(A[D]_{1,d}\)（VIII, 第 11 页）简记为 A{[}D{]}。每个元素 \(P \neq 0\)（A{[}D{]} 的）都可以唯一地写成 \(\sum_{k=0}^{m} a_k D^k\)，其中 \(a_k \in A\) 且 \(a_m \neq 0\)；我们称 m 为 P 的次数，\(a_m\) 为它的首项系数。

\begin{enumerate}
\def\labelenumi{\alph{enumi})}
\item
  设 b 是 A{[}D{]} 的一个双边理想，n 是 b 的非零元素的次数的最小值。证明由 0 与 b 中次数为 n 的元素的首项系数组成的集是 A 的一个双边理想。
\item
  设环 A 是拟单的，A 是一个 Q-代数，且导子 d 不是内导子（III, §10, No.~6, 第 560 页）。证明环 A{[}D{]} 是拟单的。
\item
  设 A 是一个特征为零的交换域，且 \(d \neq 0\)。证明 A{[}D{]} 的每个（左或右）理想都是单生成的（在这样的理想中考虑一个次数最小的元素）。由此推出 A{[}D{]} 是左且右诺特的、拟单的、无零因子的，但不包含任何极小右理想或极小左理想。
\end{enumerate}

\begin{enumerate}
\def\labelenumi{\arabic{enumi})}
\setcounter{enumi}{12}
\tightlist
\item
  设 K 是一个特征为零的交换域，A 是环 K{[}T{]}。设 d 是导子 \(\frac{d}{dT}\)。证明环 A{[}D{]}（习题 12）是拟单的。（确定导子 ad(D) 与 ad(T)，并由此推出每个非零双边理想都包含一个可逆元素。）
\end{enumerate}

¶ 14) 设 K 是一个交换域。设 M 是由两个元素 x 与 y 生成的自由幺半群，A 是 K 上幺半群 M 的代数。设 V 是 K 上的一个以 \((e_{n})_{n\geqslant1}\) 为基的向量空间。

\begin{enumerate}
\def\labelenumi{\alph{enumi})}
\item
  通过令 \(xe_1 = 0\)、\(xe_p = e_{p-1}\)（对 \(p > 1\)）及 \(ye_p = e_{2p}\)（对 \(p \geqslant 1\)），我们在 V 上定义一个 A-模结构。证明 V 是一个单 A-模。
\item
  证明 A-模 V 是忠实的。（设 \(z \in \mathbf{M}\)。对 \(n \in \mathbf{N}\)，把整数 \(s\)（使 \(e_s = ze_n\)）记作 \(\mathrm{P}_z(n)\)。注意我们有 \(\mathrm{P}_{zt} = \mathrm{P}_z \circ \mathrm{P}_t\)（对 M 中的 \(z, t\)）。对 \(z\) 与 \(t\) 的长度用归纳法由此推出：若 \(z \neq t\)，则由整数 \(n\)（使 \(\mathrm{P}_z(n) = \mathrm{P}_t(n)\)）组成的集是有限的。）
\item
  同样，对每个整数 r \textgreater{} 2，证明通过令 \(xe_{1} = 0\)、\(xe_{p} = e_{p-1}\)（对 p \textgreater{} 1）及 \(ye_{p} = e_{r^{p}}\)（对 \(p \geqslant 1\)），我们在 V 上定义了一个单且忠实的 A-模结构。证明这些结构两两不同构。
\end{enumerate}

\begin{enumerate}
\def\labelenumi{\arabic{enumi})}
\setcounter{enumi}{14}
\tightlist
\item
  设 A 是一个环。若对 A 的任意双边理想 a 与 b，关系 \(ab \subset p\) 蕴含 \(a \subset p\) 或 \(b \subset p\)，则称 A 的双边理想 p 为素理想。
\end{enumerate}

\begin{enumerate}
\def\labelenumi{\alph{enumi})}
\item
  证明当 A 交换时，这个定义与 I, §9, No.~3, 第 117 页给出的定义一致。
\item
  一个双边理想 \(\mathfrak{p}\)（\(A\) 的）是素的，当且仅当对所有元素 \(a\) 与 \(b\)（\(A - \mathfrak{p}\) 的），存在 \(x \in A\)，使得 \(axb \notin \mathfrak{p}\)。
\end{enumerate}

\begin{enumerate}
\def\labelenumi{\arabic{enumi})}
\setcounter{enumi}{15}
\tightlist
\item
  设 \(\mathrm{A}\) 是一个环。称双边理想 \(\mathfrak{p}\)（\(\mathrm{A}\) 的）为本原理想，如果商环 \(\mathrm{A} / \mathfrak{p}\) 是本原的。
\end{enumerate}

\begin{enumerate}
\def\labelenumi{\alph{enumi})}
\item
  证明本原理想就是单 A-模的零化子。特别地，交换环的本原理想就是它的极大理想。
\item
  极大双边理想是本原的；本原理想是素的（习题 15）。
\item
  A 的元素 x 是可逆的，当且仅当它在 A 对每个本原理想的每个商中的像都是可逆的（证明若情形如此，则 x 不含于任何极大左理想）。
\item
  设 \(\mathfrak{a}\) 是 A 的一个左理想。证明若有 \(\mathfrak{a} + \mathfrak{p} = \mathrm{A}\)（对 A 的每个本原理想 \(\mathfrak{p}\)），则我们有 \(\mathfrak{a} = \mathrm{A}\)。
\item
  设 M 是一个有限生成的 A-模，满足对 A 的每个本原理想 p 有 \(pM = M\)。证明 M 为零（对 M 的生成元的最小个数用归纳法，并利用 d)）。
\item
  设 M 是一个诺特 A-模。证明诸子模 aM（其中 a 取遍本原理想的有限乘积组成的集）的交约化为 0（利用 e)）。
\end{enumerate}

